\reportsection{10}{Cowork}{Cowork: Impersonating a VM Service, and Parking It}

\unithead{Mechanism}

\reppar{Cowork is the suite's largest unit and the only one with a runtime component of its own. \code{patch\_cowork\_linux()} in \code{scripts/patches/cowork.sh} runs a single embedded node heredoc against the minified \code{index.js}, version-guarded on the literal \code{vmClient (TypeScript)}, and lands roughly a dozen numbered patches that \textbf{reroute the Windows VM client to Linux}: the startVM support gate is widened past the yukonSilver feature check (FATAL on miss), the support evaluator is flipped so the renderer un-grays the tab, the multi-GB rootfs download is held disabled, the Windows named-pipe string becomes a ternary pointing at \code{\$XDG\_RUNTIME\_DIR/cowork-vm-service.sock}, empty \code{linux:\{x64:[],arm64:[]\}} manifest arrays exploit \code{[].every()} vacuous truth, and the keystone patch forks a daemon from \code{app.asar.unpacked} on \code{ENOENT}/\code{ECONNREFUSED} with a 10-second respawn cooldown, plus a quit handler that SIGTERMs it. The other half is \code{scripts/cowork-vm-service.js}, a 2,766-line daemon that \textbf{impersonates the Windows VM service} over that Unix socket, speaking the same 4-byte length-prefixed JSON protocol as the named pipe: a \code{VMManager} over Host, Bwrap, and KVM backends, bubblewrap by default and KVM opt-in via \code{COWORK\_VM\_BACKEND=kvm}, unpacked because \code{child\_process.fork} cannot execute inside an asar.}

\reppar{The same file carries two smaller passengers, the asar-path guards, addressing a shared root cause: all four legacy launchers redundantly appended the \code{app.asar} path to Electron's argv, and Electron's ASAR virtual-filesystem shim reports archives as directories, so the app dispatched its own bundle to Cowork as a folder drop. \code{patch\_asar\_path\_filter} rewrites the directory-check helper, capturing function, parameter, and fs variable dynamically:}

\begin{codeblock}
/function\s+([\w$]+)\s*\(\s*([\w$]+)\s*\)\s*\{\s*try\s*\{\s*
  return\s+([\w$]+)\.statSync\(\s*\2\s*\)\.isDirectory\(\)/
\end{codeblock}

\reppar{\code{patch\_asar\_argv\_file\_drop\_guard} covers the second-instance argv collector, which had no equivalent guard of its own: it injects the same \code{.asar}-suffix check ahead of an \code{existsSync} call, guarded by a uniqueness assertion and a TSV verification marker so a future anchor rot fails loudly instead of silently.}

\unithead{Origin}

\reppar{The reroute's origin is a same-day pivot. A \code{@ant/claude-swift} stub by @chukfinley landed via \#198 (Cowork stub PR) on 2026-02-16, shipping a documented KNOWN ISSUE (spawned-process stdout never fired, so Cowork loaded but showed nothing); hours later commit \code{25e7932} deleted the stub and introduced the daemon-plus-patch architecture that still stands. Upstream Cowork was macOS/Windows-only, so without the unit the mode was dead on Linux. @RayCharlizard is the subsystem owner, with a dedicated learning page, \code{docs/learnings/cowork-vm-daemon.md}.}

\begin{chronology}
2026-02 & \#198 \code{b8b2893}; \code{25e7932} & Origin: @chukfinley's claude-swift stub merged, then deleted hours later for the daemon-plus-\code{patch\_cowork\_linux} architecture (upstream Cowork was macOS/Windows-only). \\
2026-02 & \#259, \code{2017011} & Anchor rot: v1.1.4173 renamed the platform-gate anchor from \code{Unsupported platform} \arrow{} \code{unsupported\_platform}, crashing the app at startup; Patch 1 made a hard build error (\code{process.exit(1)}) on anchor miss, a property the June re-derivation below depends on. \\
2026-02 & \#269, \code{4929dde} & Refactor: the monolithic VM manager split into Host/Bwrap/Kvm backends, \code{COWORK\_VM\_BACKEND} override added, plus a \code{--doctor} Cowork section. \\
2026-03 & \#265, \code{d7a4606} & Guest-path fix (@leomayer): the daemon had been stripping VM guest paths from \code{--plugin-dir}/\code{--add-dir}, hiding skills and plugins; \code{translateGuestPath()} and \code{buildMountMap()} translate them instead. \\
2026-03 & \#288, PR \#300 & KVM wire-protocol series (\code{af1f2e3} through \code{473b0ba}, hardened by \code{932044d}): vsock direction/port, virtiofsd readiness, session disk plus smol-bin drive, guest RPC protocol, and a virtio-9p fallback got real VMs booting. \\
2026-03 & \#329/\#332/\#334, \code{aa6b87d}; \#337, \code{a3190c3} & Manifest fix: real Linux CDN checksums caused an infinite download-retry loop as they drifted from CDN content; superseded a day later by empty \code{linux:\{x64:[],arm64:[]\}} manifest arrays exploiting \code{[].every()} vacuous truth. \\
2026-04 & \#418, PR \#421 \code{2f6194f} (@Joost-Maker) & Identifier widening: Claude 1.3109.0 or later renamed the win32 fs var \code{e} \arrow{} \code{\$e}, breaking Patch 9 with \code{TypeError: e.existsSync is not a function}; all six extraction regexes widened from \code{(\textbackslash w+)} to \code{[\$\textbackslash w]+}. \\
2026-06 & yukonSilver, PR \#736 \code{83ea637}; \#742, \code{7327c95} & Re-derivation: upstream's yukonSilver VM refactor staled Patch 1's anchor; because Patch 1 is FATAL, main stayed red from the 1.13576.0 bump on June 17 until the June 23 re-derivation, with support-evaluator patches 1b/1c following two days later. \\
2026-04{\textendash}05 & \#383, \#622, \#632 & Symptom cluster: the ASAR virtual-filesystem shim reports \code{app.asar} as a directory, so the launchers' redundant argv path triggered a permission dialog on every launch (\#383), forced Cowork mode on every reopen (\#622), and a fatal \code{--add-dir} error in bundled Claude Code (\#632). \\
2026-05 & PR \#640, \code{6bfb296} (@aaddrick) & \code{patch\_asar\_path\_filter} lands: rewrites the directory-check helper so an \code{.asar}-suffix test short-circuits before \code{statSync().isDirectory()}. \\
2026-05 & \#668 (@MitchSchwartz); PR \#669, \code{623f1b0} & Incomplete-fix regression: the second-instance argv collector had no equivalent guard; \code{patch\_asar\_argv\_file\_drop\_guard} adds one, with a uniqueness assertion and a TSV verification marker. \\
2026-06 & PR \#700, \code{ab17b69} (@emandel82) & Root-cause fix: launchers stop passing \code{app.asar} on Electron's argv at all; the guards stay on main as defense-in-depth, load-bearing status on the deb/rpm global-Electron fallback left as unverified inference. \\
\end{chronology}

\methodbanner{\textbf{Fate under the rebase.} \bpark{} The guards are gone and the reroute is shelved, for opposite reasons. The matrix marks the guards \textbf{delete}: the official \code{statSync().isDirectory()} helpers "still exist (3 anchors, no upstream \code{.asar} guard), but the official launcher is a bare ELF symlink, no \code{app.asar} argv ever reaches them." Commit \code{d9cef9e} deleted both guard functions and the marker TSV; the new launchers exec the official binary directly, the global-Electron fallback is gone, and zero \code{endsWith(".asar")} matches remain under \code{scripts/}. The reroute itself is \textbf{park (3.1 track)}: official Cowork is a Go \code{coworkd} plus QEMU/KVM over a \code{SO\_PEERCRED} Unix socket, so the named-pipe client every anchor greps for no longer exists, and "a bwrap fallback now means impersonating that protocol, off the 3.0.0 critical path." The stack moved intact to \code{scripts/cowork-fallback/} (the reroute, a byte-identical daemon, three bats suites, and a README stating nothing there is executed or patched into any artifact); 3.0.0 ships KVM-only with doctor guidance via \code{\_check\_kvm}, \code{\_check\_vhost\_vsock}, and \code{\_check\_cowork\_stack}; the launcher keeps \code{cleanup\_orphaned\_cowork\_daemon()} to reap leftover 2.x daemons; and the \code{cowork-bwrapd} investigation against the official protocol is a 3.1 item owned by @RayCharlizard.\\[3pt]\textbf{Affects.} KVM-capable hosts; OVMF/AAVMF firmware-path sensitive (a hardcoded probe list, not distro-portable). The asar-path guards were launcher-argv defense, all Linux.}
